\documentclass[a4paper,12pt]{article}
\usepackage[utf8]{inputenc}
\usepackage[T1]{fontenc}
\usepackage[polish]{babel}
\usepackage{geometry}
\usepackage{amsmath,amssymb}
\usepackage{booktabs,longtable}
\usepackage{xcolor}
\usepackage{listings}
\usepackage[hidelinks]{hyperref}
\geometry{margin=2.5cm}

\lstdefinestyle{pythonstyle}{
  language=Python,
  basicstyle=\ttfamily\small,
  keywordstyle=\color{blue!70!black},
  commentstyle=\color{green!40!black},
  stringstyle=\color{red!60!black},
  showstringspaces=false,
  breaklines=true,
  columns=fullflexible,
  frame=single
}
\lstset{style=pythonstyle}

\title{Sprawozdanie z implementacji problemu WITI -- wariant na ocene 5.0}
\author{}
\date{}

\begin{document}
\maketitle

\section{Repozytorium}
Kod programu znajduje sie w repozytorium GitHub:
\[
\url{https://github.com/BuilderHut/komputerowo-zintegrowane-wytwarzanie/tree/main/witi}
\]

\section{Cel zadania}
Program rozwiazuje problem szeregowania z jednym procesorem, w ktorym kazde zadanie ma:
\begin{itemize}
  \item czas wykonania $p_j$,
  \item wage $w_j$,
  \item termin $d_j$.
\end{itemize}
Dla danej kolejnosci zadan liczymy calkowita kare:
\[
K = \sum_j w_j \max(C_j - d_j, 0),
\]
gdzie $C_j$ oznacza czas zakonczenia zadania $j$.

\section{Opis rozwiazania}
W pliku \texttt{witi\_solver.py} zastosowano programowanie dynamiczne po podzbiorach.
Algorytm zwraca wartosc optymalna oraz pierwsza wedlug porzadku leksykograficznego permutacje optymalna.
Nie uzywa dodatkowej tablicy pamietajacej poprzedniki albo informacje o kolejnosci.
Niech:
\begin{itemize}
  \item $P[s]$ oznacza laczny czas wykonania zadan ze zbioru $s$,
  \item $F[s]$ oznacza najmniejsza kare dla zbioru $s$, gdy zadania z tego zbioru stanowia pozostaly sufiks harmonogramu.
\end{itemize}

Niech $N$ bedzie zbiorem wszystkich zadan. Jezeli aktualny sufiks to $s$, to zadania spoza $s$ sa juz ustawione przed nim.
Czas rozpoczecia tego sufiksu wynosi:
\[
P[N] - P[s].
\]
Dla kazdego niepustego zbioru $s$ sprawdzamy, ktore zadanie $j$ moze byc pierwsze w tym sufiksie:
\[
F[s] = \min_{j \in s} \left(
  w_j \max(P[N] - P[s] + p_j - d_j, 0) + F[s \setminus \{j\}]
\right).
\]
Po wyborze pierwszego zadania sufiksu pozostaje optymalnie ulozyc zbior $s \setminus \{j\}$, dlatego rekurencja daje wynik optymalny.

\section{Rekonstrukcja pierwszej permutacji leksykograficznej}
Po wypelnieniu tablic dynamicznych odtwarzamy kolejnosc zaczynajac od zbioru wszystkich zadan.
Nie przechowujemy tablicy \texttt{last}, \texttt{parent} ani zadnych informacji o poprzednikach.
Na kazdej pozycji sprawdzamy zadania od najmniejszego numeru:
\begin{itemize}
  \item wybieramy pierwsze zadanie $j$, ktore spelnia rownanie optymalnosci,
  \item dopisujemy $j$ do permutacji,
  \item usuwamy $j$ z aktualnego zbioru,
  \item powtarzamy do wyczerpania zbioru.
\end{itemize}
Poniewaz probujemy zadania rosnaco i zawsze zachowujemy wartosc optymalna, dostajemy pierwsza leksykograficznie permutacje optymalna.
Wynik zwracany przez funkcje to minimalna kara oraz kolejnosc zadan w wersji 1-based.

\section{Zlozonosc}
Algorytm iteruje po wszystkich podzbiorach, a dla kazdego z nich sprawdza mozliwe zadania koncowe.
Zlozonosc czasowa wynosi:
\[
O(n \cdot 2^n),
\]
a zlozonosc pamieciowa:
\[
O(2^n).
\]
Odtworzenie permutacji wymaga dodatkowo $O(n^2)$ sprawdzen, co jest zdominowane przez czas wypelniania tablicy dynamicznej.
To sprawia, ze metoda jest praktyczna dla srednich instancji, ale nie dla bardzo duzych danych.

\section{Struktura projektu}
\begin{itemize}
  \item \texttt{witi\_io.py} - wczytuje kolejne instancje z pliku tekstowego,
  \item \texttt{witi\_solver.py} - zawiera programowanie dynamiczne i odbudowe pierwszej leksykograficznie permutacji optymalnej,
  \item \texttt{witi\_pd.py} - uruchamia solver i wypisuje wynik,
  \item \texttt{witi.data.txt} - zestaw danych testowych z wartosciami optymalnymi.
\end{itemize}

\section{Wyniki testowe}
Uruchomienie na calym pliku testowym daje zgodnosc z optimum dla wszystkich instancji \texttt{data.10} -- \texttt{data.20}.
Suma wynikow wynosi:
\[
\text{SUMA} = 6467.
\]

\begin{longtable}{@{}llllp{7.4cm}@{}}
\toprule
Instancja & $K$ & Optimum & Zgodnosc & Pierwsza leksykograficznie permutacja optymalna \\
\midrule
\endhead
\texttt{data.10} & 766 & 766 & tak & 6 9 2 5 1 3 4 7 8 10 \\
\texttt{data.11} & 799 & 799 & tak & 6 9 2 11 5 1 3 7 4 8 10 \\
\texttt{data.12} & 742 & 742 & tak & 6 9 2 11 5 1 3 12 7 4 8 10 \\
\texttt{data.13} & 688 & 688 & tak & 6 9 5 2 11 1 3 7 12 4 8 10 13 \\
\texttt{data.14} & 497 & 497 & tak & 6 9 5 1 2 3 11 7 12 4 8 10 13 14 \\
\texttt{data.15} & 440 & 440 & tak & 6 9 5 1 2 3 11 4 12 7 8 10 13 14 15 \\
\texttt{data.16} & 423 & 423 & tak & 6 9 5 1 2 3 7 11 4 12 8 10 13 14 15 16 \\
\texttt{data.17} & 417 & 417 & tak & 6 9 5 1 2 3 7 11 4 12 8 10 13 14 15 16 17 \\
\texttt{data.18} & 405 & 405 & tak & 6 9 5 1 2 3 7 11 12 16 17 18 4 8 10 13 14 15 \\
\texttt{data.19} & 393 & 393 & tak & 6 9 5 1 2 3 4 11 12 16 18 7 8 10 13 14 15 19 17 \\
\texttt{data.20} & 897 & 897 & tak & 6 20 9 5 1 2 3 7 11 12 18 4 8 10 13 14 16 17 19 15 \\
\bottomrule
\end{longtable}

\section{Fragment implementacji}
\lstinputlisting{witi_solver.py}

\section{Uruchomienie}
Najprostsze uruchomienie:
\begin{lstlisting}
python witi_pd.py
\end{lstlisting}
Jeden przyklad:
\begin{lstlisting}
python witi_pd.py witi.data.txt data.10
\end{lstlisting}

\end{document}
